\subsection{Architektura systemu}
\subsubsection{Moduł indeksowania dokumentów}
Diagram parsowania dokumentów i indeksowania w systemie
\subsubsection{Indeks dokumentów}
Wykorzystane indeksy - tiered index oparty o ilość cytowań, zone index dla abstractów i treści dokumentów
\subsubsection{Przetwarzanie zapytania}
Diagram sekwencji i działania
\subsubsection{Interfejs użytkownika}
Omówienie działania, zrzuty ekranu interfejsu, omówienie przykładowego workflow

\subsection{Ewaluacja jakości systemu przez użytkownika}
\subsubsection{Metryki jakości pozyskanych danych}
Omówienie precision, recall, F1
\subsubsection{Ewaluacja oceny zwracanej przez użytkownika}
Algorytm Roccio, LLM jako użytkownik końcowy przesyłający oceny, potencjalna ewaluacja przez rzeczywistego użytkownika

\subsection{Porównanie modeli w kontekście systemu wyszukiwania informacji}
\subsubsection{Metodologia}
Opis metod porównania i metryk
\subsubsection{Ewaluacja systemów bez informacji zwrotnych od użytkownika}
Porównanie modeli
\subsubsection{Ewaluacja systemów wraz z informacją zwrotną od użytkownika}
Porównanie modeli
\subsubsection{Porównanie wykorzystania zasobów i prędkości wyszukiwania}
Benchmarki pamięci i prędkości wyszukiwania end to end

\subsection{Wnioski}
Plusy i minusy naszego rozwiązania i klasycznego rozwiązania opartego o transkrypcje
\subsubsection{Porównanie modeli dla wyszukiwania bez informacji zwrotnych}
Plusy i minusy kazdego lacznie z osiagami
\subsubsection{Porównanie modeli dla wyszukiwania wraz z informacją zwrotną od użytkownika}
Plusy i minusy każdego łącznie z osiągami
\subsection{Wybór najlepszej architektury systemu do potencjalnego wykorzystania w warunkach klinicznych}
Ewaluacja wszystkich powyższych wniosków